% Generated by analysis/v3/render_taxon_context_tex.py from
% analysis/v3/taxon_context/*.tsv. Do not edit by hand.

\begin{table}[htbp]
\centering\footnotesize
\caption{Leading phyla across the $1{,}227$ core-site profiles. Values are mean shares of total reads in percent; prevalence is the fraction of profiles with at least one read. Compartment and transect-third columns are means within the stratum (thirds of $20$ sites ordered by route position). Names follow SILVA 138.2.}
\label{tab:taxa-phyla}
\begin{tabular}{lrrrrrrrr}
\toprule
Phylum & ASVs & Mean & Prev. & Surface & Shallow & Root-adj. & West & East \\
\midrule
Pseudomonadota & 88,894 & 26.9 & 100 & 29.9 & 23.8 & 27.0 & 26.3 & 28.2 \\
Bacillota & 48,693 & 22.5 & 100 & 20.8 & 19.6 & 26.2 & 18.3 & 31.5 \\
Actinomycetota & 56,581 & 18.7 & 99 & 18.0 & 20.3 & 18.1 & 22.6 & 12.7 \\
Chloroflexota & 24,409 & 8.6 & 99 & 8.5 & 10.7 & 6.9 & 10.3 & 6.0 \\
Bacteroidota & 16,657 & 7.1 & 99 & 8.2 & 5.9 & 7.3 & 6.8 & 6.0 \\
Planctomycetota & 20,808 & 3.2 & 98 & 2.3 & 3.5 & 3.8 & 3.1 & 2.9 \\
Gemmatimonadota & 12,107 & 2.7 & 98 & 2.9 & 3.5 & 1.8 & 2.7 & 2.8 \\
Acidobacteriota & 17,405 & 2.3 & 95 & 2.1 & 3.4 & 1.7 & 2.6 & 1.7 \\
Myxococcota & 15,121 & 1.7 & 98 & 1.6 & 2.0 & 1.5 & 1.7 & 1.4 \\
Deinococcota & 1,419 & 1.1 & 93 & 1.9 & 1.0 & 0.6 & 1.4 & 1.0 \\
\bottomrule
\end{tabular}
\end{table}

\begin{table}[htbp]
\centering\footnotesize
\caption{Leading genera across the $1{,}227$ core-site profiles (mean share of total reads in percent, prevalence in percent of profiles, and means by compartment and transect third).}
\label{tab:taxa-genera}
\begin{tabular}{llrrrrrrrr}
\toprule
Genus & Phylum & Mean & Prev. & Surface & Shallow & Root-adj. & West & Central & East \\
\midrule
\textit{Domibacillus} & Bacillota & 6.22 & 97 & 6.14 & 4.09 & 8.00 & 7.09 & 6.33 & 4.78 \\
\textit{Massilia} & Pseudomonadota & 3.55 & 94 & 4.85 & 2.14 & 3.66 & 3.69 & 4.58 & 1.97 \\
\textit{Bacillus} & Bacillota & 2.97 & 98 & 3.82 & 2.64 & 2.58 & 1.56 & 2.11 & 6.22 \\
\textit{Microvirga} & Pseudomonadota & 2.41 & 97 & 2.54 & 2.11 & 2.55 & 3.32 & 1.95 & 1.66 \\
\textit{Flavisolibacter} & Bacteroidota & 2.05 & 94 & 2.88 & 2.00 & 1.44 & 2.49 & 2.27 & 1.10 \\
\textit{Noviherbaspirillum} & Pseudomonadota & 1.92 & 94 & 2.62 & 1.76 & 1.50 & 2.22 & 2.10 & 1.23 \\
\textit{Streptomyces} & Actinomycetota & 1.80 & 91 & 1.43 & 1.91 & 2.01 & 2.60 & 1.71 & 0.73 \\
\textit{Herpetosiphon} & Chloroflexota & 1.52 & 92 & 1.68 & 2.07 & 0.95 & 1.45 & 2.19 & 0.72 \\
\textit{Paenibacillus} & Bacillota & 1.38 & 96 & 1.08 & 1.29 & 1.67 & 1.22 & 1.34 & 1.65 \\
\textit{Lysinibacillus} & Bacillota & 1.32 & 95 & 0.65 & 0.70 & 2.36 & 1.57 & 1.32 & 0.95 \\
\textit{Nocardioides} & Actinomycetota & 1.23 & 95 & 1.19 & 1.14 & 1.34 & 1.45 & 1.36 & 0.74 \\
\textit{Ammoniphilus} & Bacillota & 1.09 & 96 & 0.89 & 1.46 & 0.95 & 0.53 & 1.10 & 1.92 \\
\textit{Nibribacter} & Bacteroidota & 1.09 & 84 & 1.26 & 0.47 & 1.46 & 0.83 & 1.81 & 0.53 \\
\textit{Planococcus} & Bacillota & 1.03 & 82 & 1.32 & 0.60 & 1.14 & 0.74 & 0.62 & 1.99 \\
\textit{Deinococcus} & Deinococcota & 1.02 & 88 & 1.83 & 0.94 & 0.46 & 1.39 & 0.84 & 0.73 \\
\bottomrule
\end{tabular}
\end{table}

\begin{table}[htbp]
\centering\footnotesize
\caption{Genera with the strongest monotonic change along the route among the $200$ genera of the route model. $\rho$ is the Spearman correlation of the site-level CLR value (campaign-by-compartment means removed) with route position; $q$ is Benjamini--Hochberg over $200$ tests (77 genera decreased and 47 increased eastward at $q<0.05$). $\Delta$ is the east-minus-west difference in mean CLR between the transect thirds with a delete-one-site jackknife $95\,\%$ interval; the last two columns give the mean share of total reads in percent in the western and eastern thirds.}
\label{tab:taxa-replacement}
\begin{tabular}{llrrrrr}
\toprule
Genus & Phylum & $\rho$ & $q$ & $\Delta$ CLR [95\,\% CI] & West & East \\
\midrule
\multicolumn{7}{l}{\emph{Strongest decreases from west to east}} \\
\textit{Cellulomonas} & Actinomycetota & -0.79 & 2.3e-12 & -2.23 [-2.86, -1.60] & 0.50 & 0.13 \\
\textit{Blastocatella} & Acidobacteriota & -0.76 & 3.8e-11 & -2.26 [-2.80, -1.71] & 0.10 & 0.01 \\
\textit{Marmoricola} & Actinomycetota & -0.70 & 5.9e-09 & -1.65 [-2.26, -1.05] & 0.20 & 0.07 \\
\textit{Ellin6055} & Pseudomonadota & -0.70 & 7.9e-09 & -1.71 [-2.39, -1.03] & 1.12 & 0.26 \\
\textit{Roseisolibacter} & Gemmatimonadota & -0.70 & 9.3e-09 & -1.89 [-2.49, -1.28] & 0.22 & 0.09 \\
\textit{Steroidobacter} & Pseudomonadota & -0.69 & 1.2e-08 & -1.61 [-2.10, -1.12] & 0.29 & 0.11 \\
\textit{Limnobacter} & Pseudomonadota & -0.69 & 1.5e-08 & -2.83 [-3.50, -2.17] & 0.22 & 0.03 \\
\textit{Aridibacter} & Acidobacteriota & -0.69 & 1.5e-08 & -1.86 [-2.44, -1.29] & 0.08 & 0.02 \\
\textit{Sphingomonas} & Pseudomonadota & -0.67 & 5.5e-08 & -1.35 [-1.87, -0.84] & 0.21 & 0.14 \\
\textit{Paraconexibacter} & Actinomycetota & -0.67 & 5.5e-08 & -1.84 [-2.34, -1.34] & 0.05 & 0.03 \\
\midrule
\multicolumn{7}{l}{\emph{Strongest increases from west to east}} \\
\textit{Halalkalibacter} & Bacillota & 0.85 & 8.4e-16 & 4.91 [3.83, 5.99] & 0.03 & 0.83 \\
\textit{Polygonibacillus} & Bacillota & 0.81 & 6.9e-13 & 4.55 [3.59, 5.51] & 0.00 & 0.20 \\
\textit{Ammoniphilus} & Bacillota & 0.80 & 9.5e-13 & 1.89 [1.48, 2.31] & 0.53 & 1.92 \\
\textit{Bacillus} & Bacillota & 0.77 & 2.5e-11 & 1.62 [1.18, 2.07] & 1.56 & 6.22 \\
\textit{Pseudalkalibacillus} & Bacillota & 0.74 & 4.1e-10 & 2.89 [2.28, 3.50] & 0.01 & 0.05 \\
\textit{Aquibacillus} & Bacillota & 0.72 & 1.9e-09 & 4.75 [3.45, 6.04] & 0.01 & 0.69 \\
\textit{Plastorhodobacter} & Pseudomonadota & 0.72 & 3.0e-09 & 3.60 [2.74, 4.46] & 0.02 & 0.10 \\
\textit{Rhodocista} & Pseudomonadota & 0.71 & 5.9e-09 & 3.01 [2.12, 3.90] & 0.03 & 0.20 \\
\textit{Escherichia-Shigella} & Pseudomonadota & 0.70 & 7.9e-09 & 3.27 [2.36, 4.18] & 0.01 & 0.84 \\
\textit{Sediminibacillus} & Bacillota & 0.70 & 8.7e-09 & 4.59 [3.46, 5.71] & 0.00 & 0.52 \\
\bottomrule
\end{tabular}
\end{table}

\begin{table}[htbp]
\centering\footnotesize
\caption{Direction of the environmental gradients along the route: Spearman correlation of site means with route position and the means of the western and eastern transect thirds.}
\label{tab:route-gradients}
\begin{tabular}{lrrrr}
\toprule
Variable & Sites & $\rho$ & West third & East third \\
\midrule
49-month mean air temperature & 60 & 0.98 & 27.89 & 30.07 \\
49-month mean monthly rainfall & 60 & 0.55 & 2.16 & 3.67 \\
49-month mean relative humidity & 60 & 0.92 & 24.72 & 31.82 \\
archived-soil pH (admitted measurements, site mean) & 60 & 0.82 & 7.85 & 8.23 \\
\bottomrule
\end{tabular}
\end{table}

\begin{table}[htbp]
\centering\footnotesize
\caption{Genus-set overlap after subsampling every profile $20$ times to $8{,}000$ reads. A genus counts as detected in a set when it occurs in at least $10\,\%$ of subsampled profiles. ``Top 50 A in B'' is the number of the $50$ most abundant genera of set A that are detected in set B. The Atacama pit profiles (PRJEB39249) use V4 primers and intracellular DNA; the comparison is presence-based and does not merge feature tables.}
\label{tab:genus-overlap}
\begin{tabular}{p{3.6cm}p{3.6cm}rrrrrrr}
\toprule
Set A & Set B & Profiles A & Profiles B & Genera A & Genera B & Shared & Jaccard & Top 50 A in B \\
\midrule
Empty Quarter west third & Empty Quarter east third & 469 & 309 & 368 & 317 & 267 & 0.639 & 50 \\
Empty Quarter west third & Empty Quarter central third & 469 & 426 & 368 & 364 & 331 & 0.825 & 50 \\
Empty Quarter central third & Empty Quarter east third & 426 & 309 & 364 & 317 & 275 & 0.677 & 50 \\
Empty Quarter all core profiles & Atacama pit all depths (PRJEB39249) & 1204 & 58 & 364 & 76 & 43 & 0.108 & 17 \\
Empty Quarter surface & Atacama pit 2.5-10 cm (PRJEB39249) & 368 & 4 & 336 & 61 & 33 & 0.091 & 15 \\
\bottomrule
\end{tabular}
\end{table}

\begin{table}[htbp]
\centering\footnotesize
\caption{Predicted pathway classes. Share is the class total of mean relative pathway abundance across the $1{,}227$ core-site profiles (keyword grouping of MetaCyc descriptions; rules archived in \texttt{analysis/v3/taxon\_context/run\_manifest.json}). The supported counts summarize the route-correlation family ($200$ tests) and the compartment-contrast family ($600$ tests) of Section~S8.1 by class, with the number of positive (eastward increase or first-named compartment higher) and negative estimates.}
\label{tab:pathway-classes}
\begin{tabular}{lrrrrrr}
\toprule
Class & Pathways & Share (\%) & Route (+/$-$) & Compartment (+/$-$) & & \\
\midrule
Biosynthesis & 231 & 66.8 & 63 (54/9) & 208 (86/122) & & \\
Energy and central metabolism & 55 & 17.2 & 13 (6/7) & 28 (14/14) & & \\
Degradation or utilization & 135 & 10.7 & 10 (1/9) & 21 (12/9) & & \\
Other & 41 & 5.3 & 6 (5/1) & 13 (5/8) & & \\
\bottomrule
\end{tabular}
\end{table}

\begin{table}[htbp]
\centering\footnotesize
\caption{The ten most abundant predicted MetaCyc pathways and the carbon-fixation pathways present in the PICRUSt2 pathway set (mean share of predicted pathway abundance in percent; rank among $462$ pathways).}
\label{tab:pathway-dominance}
\begin{tabular}{rllr}
\toprule
Rank & Pathway & Description & Share (\%) \\
\midrule
1 & PWY-3781 & aerobic respiration I (cytochrome c) & 1.70 \\
2 & PWY-7111 & pyruvate fermentation to isobutanol (engineered) & 1.03 \\
3 & PWY-5101 & L-isoleucine biosynthesis II & 0.93 \\
4 & ILEUSYN-PWY & L-isoleucine biosynthesis I (from threonine) & 0.91 \\
5 & VALSYN-PWY & L-valine biosynthesis & 0.91 \\
6 & BRANCHED-CHAIN-AA-SYN-PWY & superpathway of branched amino acid biosynthesis & 0.80 \\
7 & PWY-5973 & cis-vaccenate biosynthesis & 0.76 \\
8 & PWY-7663 & gondoate biosynthesis (anaerobic) & 0.74 \\
9 & FASYN-ELONG-PWY & fatty acid elongation -- saturated & 0.74 \\
10 & TCA & TCA cycle I (prokaryotic) & 0.73 \\
\midrule
\multicolumn{4}{l}{\emph{Carbon-fixation pathways}} \\
44 & CALVIN-PWY & Calvin-Benson-Bassham cycle & 0.59 \\
122 & P23-PWY & reductive TCA cycle I & 0.41 \\
284 & CODH-PWY & reductive acetyl coenzyme A pathway & 0.04 \\
304 & PWY-5392 & reductive TCA cycle II & 0.02 \\
360 & PWY-5743 & 3-hydroxypropanoate cycle & 0.00 \\
362 & PWY-5744 & glyoxylate assimilation & 0.00 \\
460 & PWY-5789 & 3-hydroxypropanoate/4-hydroxybutanate cycle & 0.00 \\
\bottomrule
\end{tabular}
\end{table}
